\documentclass{situationist-preprint}

\title{The Argument From Hallucination Answered}
\msid{MND-16-b}
\msfamily{MND-16}
\msdiscipline{Philosophy of Mind}
\mstype{Companion manuscript}
\mskeywords{argument from hallucination; naïve realism; sense data; phenomenal principle; disjunctivism}

\begin{document}
\maketitle

\begin{abstract}
The argument from hallucination has two forms. The causal form infers from a shared proximal cause that perception and hallucination are of one kind; MND-16-a answers it by rejecting the principle of common kind. The phenomenological form infers from the presence of a sensory quality in hallucination that something is presented which has it, a sense datum, and then spreads the conclusion to perception. We answer the second form. Its first step relies on the phenomenal principle, which fails for hallucination, understood as sensory imagery; its spreading step relies on the common-kind principle in another guise. Naïve realism survives both forms, and hallucinations are experiences of a different kind from veridical perceptions.
\end{abstract}

\section{Two arguments}

The \emph{causal} argument from hallucination runs: a hallucination can have the same proximal cause as a veridical perception; same proximal cause, same fundamental kind; a hallucination has no external object as a constituent; so neither does a perception. MND-16-a answers it by showing that the common-kind premise (A2) is false for relational events (Lemma~4.1 and Principle~P3 there).

The \emph{phenomenological} argument has a different first step \citep{Ayer1940,Robinson1994}:
\begin{enumerate}
\item[(H1)] In a hallucination of a yellow lemon, it sensibly appears to the subject that something is yellow.
\item[(H2)] \emph{Phenomenal principle}: if it sensibly appears to a subject that something has a sensible quality, there is something of which she is aware that has it.
\item[(H3)] What she is aware of is not an external object, so it is a sense datum.
\item[(H4)] \emph{Spreading}: since a hallucination can be subjectively indistinguishable from a perception, the perception also presents a sense datum.
\end{enumerate}
If sound, it refutes naïve realism directly, whatever is said about common kinds.

\section{The answer}

\begin{proposition}\label{p:h2}
The phenomenal principle (H2) is false of hallucination.
\end{proposition}

\begin{proof}
By Assumption~5.1 of MND-16-a, a hallucination is a sensory-imagistic experience: it is of the same broad kind as vivid visualizing and dreaming. Visualizing a yellow lemon involves its sensibly appearing to one that something is yellow, in the imagistic way; it does not involve awareness of a yellow particular. If it did, imagining a ten-headed monster would bring an inner ten-headed object into existence, and the lemon one merely visualizes would be an object located somewhere. Imagery represents its objects without presenting them \citep{Martin2002}. The phenomenal principle therefore fails for imagery, and so for hallucination.
\end{proof}

This does not concede representationalism about perception. On the view of MND-16-a, the veridical perception satisfies (H2) because its object is a constituent of it (Definition~3.2 there); the hallucination fails (H2) because it is imagery. (H2) holds of one kind of experience and not of the other, which is what the claim that they differ in kind leads one to expect.

\begin{proposition}\label{p:h4}
The spreading step (H4) relies on the principle of common kind.
\end{proposition}

\begin{proof}
Indiscriminability establishes only that the subject cannot tell the two experiences apart by reflection. To infer that whatever is true of the hallucination's nature is true of the perception's, one needs a premise that indiscriminable experiences, or experiences with the same proximal cause, are of the same fundamental kind. The first is refuted by the existence of indiscriminable items of different kinds \citep{Williamson1990}; the second is (A2), refuted by Lemma~4.1 of MND-16-a.
\end{proof}

Indiscriminability is explained instead by what the two experiences share: by Lemma~5.2 of MND-16-a, they produce the same seeming that there is a yellow lemon there, and the subject's introspective judgment proceeds by how things seem. Theorem~5.3 of MND-16-a thus covers the phenomenological argument as well as the causal one.

\begin{objection}[Hallucinations are not like imagery]
Hallucinations are often taken for perceptions, while imagery never is. They cannot be of the same kind.
\end{objection}

\begin{reply}
Dreams are taken for perceptions by the dreamer, and they are imagistic. Whether imagery is taken for perception depends on its vividness and on the subject's capacity to check it against other sources, not on its fundamental kind \citep{Fish2009}. That a hallucination is taken for a perception is explained by the common seeming, which is compatible with seemings being shared by experiences that differ in kind.
\end{reply}

\begin{objection}[Spreading without common kind]
\citet{Robinson1994} argues that the brain state that suffices for the hallucination is present in the perception, so the hallucinatory experience is present there too.
\end{objection}

\begin{reply}
The brain state suffices for the seeming and, absent an object, for an imagistic experience. When the object is present, the same brain state is part of a relational event of a different kind, and nothing requires that it also produce an extra imagistic experience alongside it \citep{Martin2004}. The inference from local sufficiency to presence assumes that what a local state produces does not depend on its relations, which is Principle~P3 of MND-16-a denied.
\end{reply}

\section{Conclusion}

Both forms of the argument from hallucination rest on a premise that fails. The causal form needs the common-kind principle; the phenomenological form needs the phenomenal principle for hallucination, which fails because hallucination is imagery, and then the common-kind principle again to spread its conclusion. Naïve realism is true, and hallucinations are experiences of a different kind.

\begin{thebibliography}{9}\small
\bibitem[Ayer(1940)]{Ayer1940} Ayer, A.~J. (1940). \emph{The Foundations of Empirical Knowledge}. Macmillan.
\bibitem[Fish(2009)]{Fish2009} Fish, W. (2009). \emph{Perception, Hallucination, and Illusion}. Oxford University Press.
\bibitem[Martin(2002)]{Martin2002} Martin, M.~G.~F. (2002). The transparency of experience. \emph{Mind and Language}, 17(4), 376--425.
\bibitem[Martin(2004)]{Martin2004} Martin, M.~G.~F. (2004). The limits of self-awareness. \emph{Philosophical Studies}, 120, 37--89.
\bibitem[Robinson(1994)]{Robinson1994} Robinson, H. (1994). \emph{Perception}. Routledge.
\bibitem[Williamson(1990)]{Williamson1990} Williamson, T. (1990). \emph{Identity and Discrimination}. Blackwell.
\end{thebibliography}

\end{document}
